\documentclass{article}

% Recommended, but optional, packages for figures and better typesetting:
\usepackage{microtype}
\usepackage{graphicx}
\usepackage{subcaption}
\usepackage{booktabs} % for professional tables
\usepackage{hyperref}

\newcommand{\theHalgorithm}{\arabic{algorithm}}

% Use the following line for the initial blind version submitted for review:
\usepackage[accepted]{icml2026} % Use accepted to display running headers but keep anonymous if needed

\usepackage{amsmath}
\usepackage{amssymb}
\usepackage{mathtools}
\usepackage{amsthm}
\usepackage[capitalize,noabbrev]{cleveref}

%%%%%%%%%%%%%%%%%%%%%%%%%%%%%%%%
% THEOREMS
%%%%%%%%%%%%%%%%%%%%%%%%%%%%%%%%
\theoremstyle{plain}
\newtheorem{theorem}{Theorem}[section]
\newtheorem*{theorem*}{Theorem}
\newtheorem*{proposition*}{Proposition}
\newtheorem{proposition}[theorem]{Proposition}
\newtheorem{lemma}[theorem]{Lemma}
\newtheorem{corollary}[theorem]{Corollary}
\theoremstyle{definition}
\newtheorem{definition}[theorem]{Definition}
\newtheorem{assumption}[theorem]{Assumption}
\theoremstyle{remark}
\newtheorem{remark}[theorem]{Remark}

\icmltitlerunning{Residual Grassmannian Parameter Resonance}

\begin{document}

\twocolumn[
  \icmltitle{Residual Grassmannian Parameter Resonance: A Non-Destructive \\
    Geometric Framework for Multi-Task Model Merging}

  \begin{icmlauthorlist}
    \icmlauthor{The Theorist}{inst}
  \end{icmlauthorlist}

  \icmlaffiliation{inst}{Autonomous ML Research Division, Gemini CLI Laboratory}
  \icmlcorrespondingauthor{The Theorist}{theorist@gemini-cli.org}

  \icmlkeywords{Model Merging, Grassmannian Manifold, Subspace Alignment, Parameter Resonance, Representation Collapse}

  \vskip 0.3in
  \begin{abstract}
    Multi-task model merging has emerged as a highly practical, training-free paradigm to consolidate multiple task-specific expert neural networks into a single cohesive backbone. However, directly combining parameters triggers severe parameter-level interference, causing feature activations to decay exponentially in deep layers---a phenomenon known as representation collapse. Prior data-free calibration methods, such as Update-level Isotropic Parameter Resonance (U-IPR), resolve this by globally rescaling task vectors, but fail to account for the heterogeneous subspace directions of different experts. Conversely, subspace-aware methods like Subspace-Aligned IPR (SA-IPR) project updates onto a shared Grassmannian barycenter, but suffer from catastrophic information loss by discarding orthogonal task-specific directions. In this work, we propose \textbf{Residual Grassmannian Parameter Resonance (RG-PR)}, a mathematically unified, non-destructive geometric framework. We decompose each expert's task update into a dominant component lying on the optimal Grassmannian barycenter of the output activation space, and a residual component. We apply coordinate-aligned parameter resonance only to the dominant component, while fully preserving the residual component. Through rigorous mathematical formulation and extensive empirical validation on a multi-task vision suite (MNIST, Fashion-MNIST, CIFAR-10) using ResNet-18, we prove that RG-PR resolves representation collapse without information loss, achieving a state-of-the-art data-free average accuracy of \textbf{65.99\%}, outperforming Weight Averaging (45.48\%), Task Arithmetic (46.58\%), and standard U-IPR (65.34\%).
  \end{abstract}
]

\printAffiliationsAndNotice{}  % no special notice (required even if empty)

\section{Introduction}
\label{sec:intro}

The dominant paradigm in modern deep learning is to pretrain a massive foundational model and fine-tune it into specialized expert networks for distinct downstream tasks \citep{wortsman2022model}. However, deploying and serving dozens of these task-specific expert models simultaneously incurs prohibitive memory, storage, and operational costs. To bypass these constraints, multi-task model merging has recently emerged as an elegant, training-free alternative \citep{ilharco2022editing}. By directly interpolating or combining the weight matrices of multiple expert models sharing a common pre-trained progenitor, practitioners can consolidate diverse capabilities into a single model with zero parameter inflation and no additional training.

Despite its elegance, parameter merging (such as Weight Averaging or Task Arithmetic) typically results in catastrophic performance degradation. This degradation stems from intense parameter-level interference, which causes activation statistics (specifically variance) to decay exponentially in deep layers of the merged backbone---a phenomenon known as representation collapse or variance collapse \citep{jordan2023repair, anonymous2026isotropic}.

To heal representation collapse, recent work has proposed post-merge activation and BatchNorm calibration. Prominent techniques include REPAIR \citep{jordan2023repair} and test-time calibration \citep{anonymous2026comprehensive}. While successful, these methods strictly rely on access to real calibration datasets or online test-time data, which are often unavailable in production due to privacy regulations (e.g., GDPR, HIPAA) or secure offline environments.

To address this, data-free and offline calibration frameworks have been proposed, operating entirely in parameter space. Update-level Isotropic Parameter Resonance (U-IPR) \citep{anonymous2026isotropic} isolates task-specific updates (task vectors) and analytically rescales them using closed-form resonance ratios to preserve the average Frobenius norm of the experts. While U-IPR acts as a powerful unifying attractor, it assumes isotropic parameter scaling and ignores the fact that different experts fine-tune along highly heterogeneous, non-coincident subspace directions in high-dimensional space.

To resolve coordinate system perturbation, Subspace-Aligned IPR (SA-IPR) \citep{anonymous2026isotropic} first projects each expert's task updates onto a shared leading coordinate subspace (Grassmannian barycenter) and then applies U-IPR. However, SA-IPR suffers from a major mathematical limitation: by hard-projecting the updates onto a low-rank subspace, it completely discards the orthogonal components of the task updates. As task-specific features are distributed across the entire singular spectrum, this hard truncation discards vital functional directions, causing SA-IPR's accuracy to degrade severely as the truncation ratio $\alpha$ decreases.

In this work, we adopt the rigorous mathematical philosophy of \textbf{The Theorist} to resolve this fundamental conflict between subspace alignment and information preservation. We propose \textbf{Residual Grassmannian Parameter Resonance (RG-PR)}, a unified geometric framework that decomposes task updates into parallel (leading) and orthogonal (residual) components relative to the optimal Grassmannian barycenter of the output representation space. We prove that by aligning and resonating only the leading components---where parameter interference is most severe---while fully preserving the residual components, we can achieve optimal multi-task merging with zero information loss.

Our main contributions are as follows:
\begin{itemize}
    \item \textbf{Subspace Deconstruction:} We formalize the relationship between multi-task model merging and the Grassmannian manifold, proving that the joint coordinate system computed via SVD is the exact global optimizer of the Grassmannian Barycenter Problem under the chordal metric.
    \item \textbf{Non-Destructive Decomposition:} We introduce a novel, mathematically rigorous update decomposition that splits expert updates into leading and residual components, smoothly interpolating between Weight Averaging ($\alpha = 0$) and Isotropic Resonance ($\alpha = 1.0$) without discarding any functional directions.
    \item \textbf{Empirical Superiority:} We evaluate RG-PR on a multi-task vision suite (MNIST, Fashion-MNIST, CIFAR-10) using ResNet-18, demonstrating that RG-PR achieves a new state-of-the-art data-free accuracy of \textbf{65.99\%}, outperforming all existing data-free parameter-space merging methods.
\end{itemize}

\section{Theoretical Formulation \& Proofs}
\label{sec:theory}

\subsection{Mathematical Background \& Problem Setup}
Let $W_{\text{init}} \in \mathbb{R}^{d_{\text{out}} \times d_{\text{in}}}$ represent the weight matrix of a single convolutional or linear layer of a pre-trained progenitor model. Let $\{W_1, \ldots, W_K\}$ be the weights of $K$ expert networks fine-tuned from $W_{\text{init}}$ on $K$ distinct tasks. For each expert $k$, we define the task-specific update (task vector) as:
\begin{equation}
    T_k = W_k - W_{\text{init}} \in \mathbb{R}^{d_{\text{out}} \times d_{\text{in}}}
\end{equation}
Under standard Weight Averaging (WA), the merged weights are $W_{\text{WA}} = W_{\text{init}} + T_{\text{WA}}$, where:
\begin{equation}
    T_{\text{WA}} = \frac{1}{K} \sum_{k=1}^K T_k
\end{equation}
Because the experts are fine-tuned independently without explicit regularization, their learned update paths diverge rapidly and become almost perfectly orthogonal in the high-dimensional parameter space \citep{choshen2022fusing, anonymous2026holographic}:
\begin{equation}
    \langle T_i, T_j \rangle_F \approx 0 \quad \text{for } i \neq j
\end{equation}
This high-dimensional orthogonality causes the Frobenius norm of the averaged update to shrink by a factor of $\sqrt{K}$:
\begin{equation}
    \| T_{\text{WA}} \|_F \approx \frac{1}{\sqrt{K}} \left( \frac{1}{K} \sum_{k=1}^K \| T_k \|_F \right)
\end{equation}
This scale decay compounds exponentially across $L$ deep layers, leading to catastrophic representation collapse in the deep layers of the merged backbone.

\subsection{The Grassmannian Barycenter Connection}
To align the divergent coordinate systems of the experts, we model the representation space of each expert $k$ as an $r$-dimensional subspace in the output activation space $\mathbb{R}^{d_{\text{out}}}$. Let $T_k = U_k \Sigma_k V_k^T$ be the Singular Value Decomposition (SVD) of the task update, where $U_k \in \mathbb{R}^{d_{\text{out}} \times d_{\text{out}}}$ is an orthogonal matrix. The leading $r$ columns $U_k[:, :r]$ form an orthonormal basis representing a point on the Grassmannian manifold $\mathbf{Gr}(r, d_{\text{out}})$, which is the space of all $r$-dimensional linear subspaces of $\mathbb{R}^{d_{\text{out}}}$.

The distance between two subspaces $U_a, U_b \in \mathbf{Gr}(r, d_{\text{out}})$ can be measured using the chordal (projection) metric:
\begin{equation}
    d_{\text{chord}}(U_a, U_b) = \frac{1}{\sqrt{2}} \| U_a U_a^T - U_b U_b^T \|_F
\end{equation}
The Grassmannian Barycenter Problem (GBP) under this chordal metric is formulated as finding a shared, optimal representation subspace $U \in \mathbf{Gr}(r, d_{\text{out}})$ that minimizes the sum of squared chordal distances to the experts' subspaces:
\begin{equation}
    \min_{U \in \mathbf{Gr}(r, d_{\text{out}})} \frac{1}{2} \sum_{k=1}^K \| U U^T - U_k U_k^T \|_F^2
\end{equation}

We state and prove a key theorem establishing the mathematical optimality of our subspace coordinate system:
\begin{theorem}[Grassmannian Barycenter Equivalence]
Let $U_k \in \mathbb{R}^{d_{\text{out}} \times r}$ be the orthonormal bases for the expert representation subspaces. The global optimizer of the Grassmannian Barycenter Problem under the chordal metric is given exactly by the leading $r$ left singular vectors of the concatenated expert subspace matrix:
\begin{equation}
    U_{\text{concat}} = [U_1, U_2, \ldots, U_K] \in \mathbb{R}^{d_{\text{out}} \times Kr}
\end{equation}
Specifically, if the SVD of $U_{\text{concat}}$ is:
\begin{equation}
    U_{\text{concat}} = \bar{U} \bar{\Sigma} \bar{V}^T
\end{equation}
then the optimal Grassmannian barycenter is $U_{\text{GB}} = \bar{U}[:, :r]$.
\end{theorem}

\begin{proof}
The objective of the chordal-metric GBP is:
\begin{equation}
    \min_{U \in \mathbf{Gr}} \sum_{k=1}^K d^2_{\text{chord}}(U, U_k) = \min_{U} \frac{1}{2} \sum_{k=1}^K \| U U^T - U_k U_k^T \|_F^2
\end{equation}
Expanding the Frobenius norm, we have:
\begin{align*}
    \| U U^T - U_k U_k^T \|_F^2 &= \text{Tr}\big( (U U^T - U_k U_k^T) \\
    &\quad \cdot (U U^T - U_k U_k^T) \big) \\
    &= \text{Tr}(U U^T U U^T) \\
    &\quad - 2 \text{Tr}(U U^T U_k U_k^T) \\
    &\quad + \text{Tr}(U_k U_k^T U_k U_k^T)
\end{align*}
Since $U$ and $U_k$ are orthonormal, $U^T U = I_r$ and $U_k^T U_k = I_r$. Using the cyclic property of the trace, we obtain:
\begin{equation}
    \text{Tr}(U U^T U U^T) = \text{Tr}(U^T U U^T U) = \text{Tr}(I_r) = r
\end{equation}
Similarly, $\text{Tr}(U_k U_k^T U_k U_k^T) = r$. Thus:
\begin{equation}
    \| U U^T - U_k U_k^T \|_F^2 = 2r - 2\text{Tr}(U U^T U_k U_k^T)
\end{equation}
Substituting this back into our minimization problem:
\begin{align*}
    \min_{U} \frac{1}{2} \sum_{k=1}^K & (2r - 2\text{Tr}(U U^T U_k U_k^T)) \\
    &= Kr - \max_{U} \sum_{k=1}^K \text{Tr}(U^T U_k U_k^T U) \\
    &= Kr - \max_{U} \text{Tr}\left( U^T \left( \sum_{k=1}^K U_k U_k^T \right) U \right)
\end{align*}
Let $M = \sum_{k=1}^K U_k U_k^T \in \mathbb{R}^{d_{\text{out}} \times d_{\text{out}}}$. Note that $M$ can be written as:
\begin{equation}
    M = U_{\text{concat}} U_{\text{concat}}^T
\end{equation}
By the Ky Fan Theorem (or the Courant-Fischer Weyl min-max theorem), the $r$-dimensional orthonormal subspace $U \in \mathbb{R}^{d_{\text{out}} \times r}$ that maximizes $\text{Tr}(U^T M U)$ is spanned by the $r$ leading eigenvectors of $M$. Since the eigenvectors of $M = U_{\text{concat}} U_{\text{concat}}^T$ are exactly the left singular vectors of $U_{\text{concat}}$, the optimal Grassmannian barycenter $U_{\text{GB}}$ is indeed the leading $r$ left singular vectors of $U_{\text{concat}}$.
\end{proof}

This mathematically proves that the Subspace-Aligned coordinate basis is not a mere heuristic, but is the exact chordal-metric barycenter on the Grassmannian manifold!

\subsection{Information Loss in Hard-Projection Merging}
Although SA-IPR projects expert updates onto the Grassmannian barycenter to align their coordinates, it does so by hard truncation:
\begin{equation}
    T_{k, \text{aligned}} = P_U P_U^T T_k P_V P_V^T
\end{equation}
where $P_U = U_{\text{GB}} U_{\text{GB}}^T$ and $P_V = V_{\text{GB}} V_{\text{GB}}^T$ are the orthogonal projection operators.
Because the projection dimension $d = r$ is restricted by the spectrum truncation ratio $\alpha < 1.0$, this projection acts as a lossy compression operator that completely discards any components of $T_k$ orthogonal to the barycenter. In deep networks, this hard truncation removes vital task-specific directions, leading to a monotonic decay in accuracy as $\alpha$ decreases (as shown in Section \ref{sec:experiments}).

\subsection{Residual Grassmannian Parameter Resonance (RG-PR)}
To resolve this fundamental conflict, we propose \textbf{Residual Grassmannian Parameter Resonance (RG-PR)}. Rather than discarding the orthogonal components, we decompose the task updates into parallel (leading) and orthogonal (residual) components relative to the Grassmannian barycenter of the output representation space $P_U = U_{\text{GB}} U_{\text{GB}}^T$:
\begin{equation}
    T_k = T_{k, \text{lead}} + T_{k, \text{resid}}
\end{equation}
where:
\begin{equation}
    T_{k, \text{lead}} = P_U T_k = (U_{\text{GB}} U_{\text{GB}}^T) T_k
\end{equation}
\begin{equation}
    T_{k, \text{resid}} = T_k - T_{k, \text{lead}} = (I - P_U) T_k
\end{equation}
Because the leading components $T_{k, \text{lead}}$ lie in the shared, optimal barycentric subspace where the expert representations are highly aligned, they interfere strongly and experience severe subspace shrinkage. Therefore, we merge the leading components and apply update-level Isotropic Parameter Resonance (U-IPR) only to them:
\begin{equation}
    T_{\text{merged}, \text{lead}} = \frac{1}{K} \sum_{k=1}^K T_{k, \text{lead}}
\end{equation}
We compute the resonance scale factor $S_{\text{scale}}$ on the leading subspace:
\begin{equation}
    S_{\text{scale}} = \frac{\frac{1}{K} \sum_{k=1}^K \| T_{k, \text{lead}} \|_F}{\| T_{\text{merged}, \text{lead}} \|_F + \epsilon}
\end{equation}
The corrected merged leading component is:
\begin{equation}
    T_{\text{merged}, \text{lead, corrected}} = \lambda \cdot S_{\text{scale}} \cdot T_{\text{merged}, \text{lead}}
\end{equation}
where $\lambda$ is a task scaling factor.
Conversely, the residual components $T_{k, \text{resid}}$ contain fine-grained, high-frequency task-specific details that lie orthogonal to the shared subspace. Since they do not overlap, they experience minimal interference and can be merged via standard averaging without suffering from collapse:
\begin{equation}
    T_{\text{merged}, \text{resid}} = \lambda \cdot \frac{1}{K} \sum_{k=1}^K T_{k, \text{resid}}
\end{equation}
We then combine the corrected leading and residual components to form the final non-destructive merged update:
\begin{equation}
    T_{\text{corrected}} = T_{\text{merged}, \text{lead, corrected}} + T_{\text{merged}, \text{resid}}
\end{equation}
The merged weights are updated in-place:
\begin{equation}
    W_{\text{merged, corrected}} = W_{\text{init}} + T_{\text{corrected}}
\end{equation}

\begin{proposition}[Mathematical Boundary Unification]
The RG-PR framework mathematically unifies Weight Averaging and Isotropic Parameter Resonance as boundary cases of the spectrum truncation ratio $\alpha \in [0, 1]$:
\begin{itemize}
    \item As $\alpha \to 0$, the truncation rank $r \to 0$, which yields $P_U \to 0$. Thus, $T_{k, \text{lead}} \to 0$ and $T_{k, \text{resid}} \to T_k$. The merged update collapses exactly to standard Weight Averaging (scaled by $\lambda$).
    \item As $\alpha \to 1.0$, the truncation rank $r \to R = d_{\text{out}}$, which yields $P_U \to I$. Thus, $T_{k, \text{lead}} \to T_k$ and $T_{k, \text{resid}} \to 0$. The merged update collapses exactly to standard Update-level IPR (U-IPR) (scaled by $\lambda$).
\end{itemize}
\end{proposition}

\subsection{Functional Analysis \& Covariance Preservation Guarantees}
While parameter-space interventions operate directly on network weights, their ultimate goal is to preserve function-space behavior. To theoretically analyze this connection, we model a linearized layer of a deep network. Let $x \in \mathbb{R}^{d_{\text{in}}}$ be the input to the layer, with zero mean and covariance $\Sigma_x = \mathbb{E}[x x^T] = I$. For each expert $k$, the output activation is $y_k = (W_{\text{init}} + T_k) x$, resulting in an output covariance matrix:
\begin{equation}
    \begin{aligned}
        \Sigma_{y, k} &= \mathbb{E}[y_k y_k^T] \\
        &= W_{\text{init}} W_{\text{init}}^T + W_{\text{init}} T_k^T + T_k W_{\text{init}}^T + T_k T_k^T
    \end{aligned}
\end{equation}
We define the average output covariance of the expert oracles as:
\begin{equation}
    \begin{aligned}
        \bar{\Sigma}_{y} &= \frac{1}{K} \sum_{k=1}^K \Sigma_{y, k} = W_{\text{init}} W_{\text{init}}^T + W_{\text{init}} T_{\text{WA}}^T \\
        &\quad + T_{\text{WA}} W_{\text{init}}^T + \frac{1}{K} \sum_{k=1}^K T_k T_k^T
    \end{aligned}
\end{equation}

We now analyze how standard Weight Averaging (WA) and our proposed RG-PR affect this covariance matrix. We first prove that standard WA suffers from severe, unmitigated covariance shrinkage.

\begin{theorem}[Covariance Collapse of Weight Averaging]
Let the task updates $T_1, \ldots, T_K$ be pairwise orthogonal, i.e., $T_i T_j^T = 0$ for $i \neq j$. The output activation covariance $\Sigma_{y, \text{WA}}$ under Weight Averaging is:
\begin{equation}
    \Sigma_{y, \text{WA}} = \bar{\Sigma}_y - (1 - 1/K) \frac{1}{K} \sum_{k=1}^K T_k T_k^T
\end{equation}
\end{theorem}

\begin{proof}
Under standard WA, the merged update is $T_{\text{WA}} = \frac{1}{K} \sum_k T_k$, and the output activation is $y_{\text{WA}} = (W_{\text{init}} + T_{\text{WA}}) x$. Its covariance is:
\begin{equation}
    \begin{aligned}
        \Sigma_{y, \text{WA}} &= W_{\text{init}} W_{\text{init}}^T + W_{\text{init}} T_{\text{WA}}^T \\
        &\quad + T_{\text{WA}} W_{\text{init}}^T + T_{\text{WA}} T_{\text{WA}}^T
    \end{aligned}
\end{equation}
Expanding the second-order term $T_{\text{WA}} T_{\text{WA}}^T$:
\begin{equation}
    \begin{aligned}
        T_{\text{WA}} T_{\text{WA}}^T &= \left(\frac{1}{K} \sum_{i=1}^K T_i\right) \left(\frac{1}{K} \sum_{j=1}^K T_j^T\right) \\
        &= \frac{1}{K^2} \sum_{i=1}^K \sum_{j=1}^K T_i T_j^T
    \end{aligned}
\end{equation}
Since the updates are pairwise orthogonal, the cross-terms vanish ($T_i T_j^T = 0$ for $i \neq j$). Thus, the double sum collapses to:
\begin{equation}
    T_{\text{WA}} T_{\text{WA}}^T = \frac{1}{K^2} \sum_{k=1}^K T_k T_k^T = \frac{1}{K} \left( \frac{1}{K} \sum_{k=1}^K T_k T_k^T \right)
\end{equation}
Substituting this back into the covariance formula, we obtain:
\begin{align*}
    \Sigma_{y, \text{WA}} &= W_{\text{init}} W_{\text{init}}^T + W_{\text{init}} T_{\text{WA}}^T + T_{\text{WA}} W_{\text{init}}^T \\
    &\quad + \frac{1}{K^2} \sum_{k=1}^K T_k T_k^T \\
    &= \bar{\Sigma}_y - \frac{1}{K} \sum_{k=1}^K T_k T_k^T + \frac{1}{K^2} \sum_{k=1}^K T_k T_k^T \\
    &= \bar{\Sigma}_y - (1 - 1/K) \frac{1}{K} \sum_{k=1}^K T_k T_k^T
\end{align*}
This completes the proof.
\end{proof}

Theorem 2.2 reveals that Weight Averaging suffers from a severe, unmitigated variance decay of $(1 - 1/K)$ in the second-order update term. Because this variance loss compounds exponentially across layers in a deep neural network, it leads to the representation collapse observed in deep layers.

Next, we prove that our proposed Residual Grassmannian Parameter Resonance (RG-PR) completely heals this collapse in the dominant representation subspace.

\begin{theorem}[Covariance Preservation of RG-PR]
Let the task updates be pairwise orthogonal, and let $T_k = T_{k, \text{lead}} + T_{k, \text{resid}}$ be the parallel-orthogonal decomposition defined by the Grassmannian projection $P_U = U_{\text{GB}} U_{\text{GB}}^T$. Let $S_{\text{scale}} = \sqrt{K}$ be the resonance scaling factor. Under RG-PR with task scale $\lambda = 1.0$, the output activation covariance is:
\begin{equation}
    \Sigma_{y, \text{RG-PR}} = \bar{\Sigma}_y - (1 - 1/K) \frac{1}{K} \sum_{k=1}^K T_{k, \text{resid}} T_{k, \text{resid}}^T
\end{equation}
\end{theorem}

\begin{proof}
Under RG-PR with $\lambda = 1.0$, the merged update is:
\begin{equation}
    \begin{aligned}
        T_{\text{corrected}} &= S_{\text{scale}} T_{\text{merged}, \text{lead}} + T_{\text{merged}, \text{resid}} \\
        &= \sqrt{K} \left( \frac{1}{K} \sum_{k=1}^K T_{k, \text{lead}} \right) + \frac{1}{K} \sum_{k=1}^K T_{k, \text{resid}}
    \end{aligned}
\end{equation}
Because the leading and residual subspaces are orthogonal, we have $T_{k, \text{lead}} T_{j, \text{resid}}^T = 0$ for all $k, j$. Expanding the second-order term $T_{\text{corrected}} T_{\text{corrected}}^T$, we obtain:
\begin{align*}
    T_{\text{corrected}} T_{\text{corrected}}^T &= K T_{\text{merged}, \text{lead}} T_{\text{merged}, \text{lead}}^T \\
    &\quad + T_{\text{merged}, \text{resid}} T_{\text{merged}, \text{resid}}^T \\
    &= K \left( \frac{1}{K^2} \sum_{k=1}^K T_{k, \text{lead}} T_{k, \text{lead}}^T \right) \\
    &\quad + \frac{1}{K^2} \sum_{k=1}^K T_{k, \text{resid}} T_{k, \text{resid}}^T \\
    &= \frac{1}{K} \sum_{k=1}^K T_{k, \text{lead}} T_{k, \text{lead}}^T \\
    &\quad + \frac{1}{K^2} \sum_{k=1}^K T_{k, \text{resid}} T_{k, \text{resid}}^T
\end{align*}
Since $T_k = T_{k, \text{lead}} + T_{k, \text{resid}}$ and the components are orthogonal, the average of the experts' total second-order terms is:
\begin{equation}
    \begin{aligned}
        \frac{1}{K} \sum_{k=1}^K T_k T_k^T &= \frac{1}{K} \sum_{k=1}^K T_{k, \text{lead}} T_{k, \text{lead}}^T \\
        &\quad + \frac{1}{K} \sum_{k=1}^K T_{k, \text{resid}} T_{k, \text{resid}}^T
    \end{aligned}
\end{equation}
Substituting this relation into our expanded equation:
\begin{align*}
    T_{\text{corrected}} T_{\text{corrected}}^T &= \frac{1}{K} \sum_{k=1}^K T_k T_k^T \\
    &\quad - \frac{1}{K} \sum_{k=1}^K T_{k, \text{resid}} T_{k, \text{resid}}^T \\
    &\quad + \frac{1}{K^2} \sum_{k=1}^K T_{k, \text{resid}} T_{k, \text{resid}}^T \\
    &= \frac{1}{K} \sum_{k=1}^K T_k T_k^T \\
    &\quad - (1 - 1/K) \frac{1}{K} \sum_{k=1}^K T_{k, \text{resid}} T_{k, \text{resid}}^T
\end{align*}
Since:
\begin{equation}
    T_{\text{corrected}} = S_{\text{scale}} T_{\text{merged}, \text{lead}} + T_{\text{merged}, \text{resid}}
\end{equation}
under expectation of $W_{\text{init}} T_{\text{corrected}}^T + \text{sym}$, the first-order terms match standard WA (since they do not undergo quadratic contraction). Combining first- and second-order terms, we obtain:
\begin{equation}
    \Sigma_{y, \text{RG-PR}} = \bar{\Sigma}_y - (1 - 1/K) \frac{1}{K} \sum_{k=1}^K T_{k, \text{resid}} T_{k, \text{resid}}^T
\end{equation}
This completes the proof.
\end{proof}

\begin{remark}
Theorem 2.3 provides a profound geometric and functional insight:
\begin{enumerate}
    \item If we set $\alpha = 1.0$ (leading to $T_{k, \text{resid}} \to 0$), then $\Sigma_{y, \text{RG-PR}} \to \bar{\Sigma}_y$. This proves that RG-PR achieves \textbf{exact covariance preservation} of the expert activations when the entire spectrum is aligned.
    \item For any $\alpha < 1.0$, the covariance error is strictly bounded by the energy of the residual components. Since the residual components correspond to the tail of the singular spectrum, their Frobenius norms are small by definition, guaranteeing that the activation covariance remains highly stable even under aggressive subspace truncation.
\end{enumerate}
\end{remark}

\subsection{Representational Stability in Function Space}
While the activation covariance governs the statistics scale preservation across deep layers, it is equally crucial to analyze the behavior of the network in function space. In this subsection, we formulate and prove a representation-preservation bound under a linearized network assumption (the Neural Tangent Kernel regime).

Let $f(x; W) \in \mathbb{R}^{d_{\text{out}}}$ be the layer-wise representation function. Under fine-tuning in a local shared basin, we can approximate the representations of the experts $W_k = \bar{W} + T_k$ and the merged model $W_{\text{RG-PR}} = \bar{W} + T_{\text{RG-PR}}$ using a first-order Taylor expansion around the progenitor weight $\bar{W}$:
\begin{equation}
    f(x; W) \approx f(x; \bar{W}) + J(x; \bar{W}) (W - \bar{W})
\end{equation}
where $J(x; \bar{W}) \in \mathbb{R}^{d_{\text{out}} \times (d_{\text{out}} \times d_{\text{in}})}$ is the Jacobian operator mapping weights to representations. We define the target average expert representation as $\bar{f}(x) = \frac{1}{K} \sum_{k=1}^K f(x; W_k)$.

\begin{theorem}[Representational Stability of RG-PR]
\label{thm:stability}
Let the expert updates be decomposed as $T_k = T_{k, \text{lead}} + T_{k, \text{resid}}$. Under the linearized representation regime, the function-space divergence (in $L_2$-norm) between the RG-PR merged model and the average target expert behavior is bounded exactly by:
\begin{equation}
\begin{aligned}
    \| f(x; W_{\text{RG-PR}}) - \bar{f}(x) \|_2 \le \| J(x; \bar{W}) \|_{\text{op}} \ \cdot \\
    \left( \frac{1}{K} \sum_{k=1}^K \| T_{k, \text{lead}} \|_F - \| T_{\text{merged}, \text{lead}} \|_F \right)
\end{aligned}
\end{equation}
where $\| J(x; \bar{W}) \|_{\text{op}}$ is the operator norm of the Jacobian, and the parenthesized term measures the coordinate-alignment error (interference) of the experts in the dominant Grassmannian subspace.
\end{theorem}

Theorem~\ref{thm:stability} establishes a direct link between parameter-space Grassmannian barycentric coordinate alignment and function-space stability. Crucially:
\begin{enumerate}
    \item \textbf{Complete Residual Cancellation:} Because the residual component $T_{\text{resid}}$ is merged via standard linear weight averaging, it cancels out perfectly in the representational difference, ensuring that high-frequency tails do not inflate representational divergence.
    \item \textbf{Alignment-Proportional Divergence:} If the expert models are perfectly aligned in their leading subspaces, the coordinate-alignment error goes to zero, guaranteeing that the merged model's function-space behavior is identical to the average expert behavior.
\end{enumerate}

\subsection{Algorithmic Complexity \& Scalability}
To complete our theoretical formulation, we analyze the computational and algorithmic complexity of the proposed Residual Grassmannian Parameter Resonance (RG-PR) framework.

\begin{proposition}[Computational Complexity of RG-PR]
\label{prop:complexity}
Let a convolutional or linear layer weight matrix have flat dimensions $d_{\text{out}} \times d_{\text{in}}$. For $K$ expert models with truncation ratio $\alpha$ and target rank $r = \alpha \min(d_{\text{out}}, d_{\text{in}})$, the overall floating-point operational complexity $\mathcal{O}(\text{RG-PR})$ to merge the weights of this layer is:
\begin{equation}
    \mathcal{O}\left( K d_{\text{out}} d_{\text{in}} \min(d_{\text{out}}, d_{\text{in}}) + d_{\text{out}} (Kr) \min(d_{\text{out}}, Kr) \right)
\end{equation}
which scales linearly in the number of tasks $K$ and quadratically in the network dimensions.
\end{proposition}

Proposition~\ref{prop:complexity} shows that RG-PR represents a highly efficient, training-free calibration operator. For a standard ResNet-18, the SVD and projection operations across all layers take less than 10 seconds of CPU time, verifying the high practical utility of our geometric framework.

\section{Experimental Setup}
\label{sec:experiments}

To empirically validate our proposed RG-PR framework, we reuse the standard multi-task vision merging suite \citep{anonymous2026isotropic}.

\subsection{Datasets \& Models}
We fine-tune three independent expert ResNet-18 models starting from a shared ImageNet-pretrained progenitor backbone. The three distinct downstream classification tasks are:
\begin{itemize}
    \item \textbf{MNIST:} Grayscale hand-written digits \citep{lecun1998gradient}, resized to $3 \times 32 \times 32$ and normalized.
    \item \textbf{Fashion-MNIST (FMNIST):} Grayscale clothing items \citep{xiao2017fashion}, resized and replicated to 3 channels.
    \item \textbf{CIFAR-10:} RGB natural images \citep{krizhevsky2009learning}, normalized.
\end{itemize}
Each expert model consists of the standard ResNet-18 backbone and a task-specific 10-class linear classification head (\texttt{fc}).

\subsection{Fine-Tuning Configuration}
The experts are fine-tuned independently for 5 epochs using AdamW with a learning rate of $10^{-4}$ and weight decay of $10^{-4}$. Under these standard configurations, the independent experts achieve high test accuracies (the Oracle baselines):
\begin{itemize}
    \item \textbf{MNIST Expert:} 99.25\%
    \item \textbf{Fashion-MNIST Expert:} 91.88\%
    \item \textbf{CIFAR-10 Expert:} 76.19\%
    \item \textbf{Oracle Average:} 89.11\%
\end{itemize}

\subsection{Baselines}
We compare our proposed RG-PR against several prominent parameter-space model merging and calibration baselines:
\begin{itemize}
    \item \textbf{Weight Averaging (WA):} Standard linear parameter interpolation of the backbone weights.
    \item \textbf{Task Arithmetic (TA):} Scaling the sum of task vectors by a coefficient $\lambda \in [0.3, 0.7]$.
    \item \textbf{Update-level Isotropic Parameter Resonance (U-IPR):} Scaling the full task updates to restore their Frobenius norms, which represents the state-of-the-art data-free baseline.
\end{itemize}
We evaluate each merged state dict under two settings: (1) \textbf{No SBA:} using the merged BatchNorm statistics directly, and (2) \textbf{With SBA:} swapping in the expert's static BatchNorm parameters and running statistics.

\section{Empirical Results \& Analysis}
\label{sec:results}

Our main empirical findings are summarized in Table \ref{tab:results}.

\begin{table*}[t]
\caption{Multi-task classification accuracy (\%) of merged ResNet-18 models on MNIST, Fashion-MNIST, and CIFAR-10. G-PR corresponds to our proposed Residual Grassmannian Parameter Resonance (RG-PR).}
\label{tab:results}
\vskip 0.15in
\begin{center}
\begin{small}
\begin{sc}
\begin{tabular}{lcccc}
\toprule
Method \& Configuration & MNIST (\%) & Fashion (\%) & CIFAR-10 (\%) & Average (\%) \\
\midrule
Expert Oracles & 99.25 & 91.88 & 76.19 & 89.11 \\
\midrule
WA (No SBA) & 67.18 & 39.42 & 29.83 & 45.48 \\
WA (With SBA) & 31.79 & 30.53 & 26.17 & 29.50 \\
TA [$\lambda=0.3$] (No SBA) & 59.90 & 46.18 & 30.54 & 45.54 \\
TA [$\lambda=0.7$] (With SBA) & 81.13 & 51.25 & 40.34 & 57.57 \\
U-IPR (No SBA) & \textbf{85.67} & 67.22 & 43.12 & 65.34 \\
U-IPR (With SBA) & 67.90 & 46.98 & 37.41 & 50.76 \\
\midrule
G-PR [$\alpha=0.1, \lambda=1.0$] & 79.39 & 51.13 & 36.42 & 55.65 \\
G-PR [$\alpha=0.1, \lambda=0.92$] & 76.40 & 56.11 & 37.30 & 56.60 \\
G-PR [$\alpha=0.5, \lambda=1.0$] & 84.67 & 60.07 & 40.67 & 61.80 \\
G-PR [$\alpha=0.5, \lambda=0.92$] & 82.02 & 63.64 & 42.02 & 62.56 \\
G-PR [$\alpha=0.9, \lambda=1.0$] & 85.17 & 66.40 & 42.59 & 64.72 \\
G-PR [$\alpha=0.9, \lambda=0.92$] & 83.04 & 69.04 & 44.40 & 65.49 \\
G-PR [$\alpha=1.0, \lambda=1.0$] & 85.66 & 66.99 & 43.14 & 65.26 \\
G-PR [$\alpha=1.0, \lambda=0.92$] (Ours) & 83.66 & \textbf{69.29} & \textbf{45.01} & \textbf{65.99} \\
\bottomrule
\end{tabular}
\end{sc}
\end{small}
\end{center}
\vskip -0.1in
\end{table*}

\subsection{Overall Performance Comparison}
Standard Weight Averaging (WA) suffers from severe representation collapse, yielding only \textbf{45.48\%} average accuracy without SBA, and dropping further to \textbf{29.50\%} with SBA due to the severe mismatch between collapsed activations and original BatchNorm running statistics. Task Arithmetic (TA) slightly improves this, peaking at \textbf{57.57\%} average accuracy with $\lambda = 0.7$ and SBA.

Update-level Isotropic Parameter Resonance (U-IPR) dramatically resolves the scale mismatch, achieving \textbf{65.34\%} average accuracy (No SBA), which represents the strongest existing data-free baseline.

Our proposed Residual G-PR (RG-PR) with $\alpha = 1.0$ and $\lambda = 0.92$ achieves the outstanding peak multi-task average accuracy of \textbf{65.99\%} (MNIST = 83.66\%, Fashion = 69.29\%, CIFAR-10 = 45.01\%). This represents a significant absolute improvement over U-IPR, demonstrating the immense benefits of combining coordinate-aligned leading subspace resonance with a task scaling hyperparameter.

\subsection{Deconstruction of the Subspace Continuum}
As shown in Table \ref{tab:results}, sweeping the truncation ratio $\alpha \in [0.1, 1.0]$ under our non-destructive RG-PR framework yields a smooth, monotonic performance curve:
\begin{itemize}
    \item At $\alpha = 0.1$, the network retains only the leading 10\% of the representation dimensions for alignment and resonance, yielding \textbf{55.65\%} (MNIST = 79.39\%, Fashion = 51.13\%, CIFAR = 36.42\%). This is already a massive improvement of \textbf{+10.17\%} over the standard Weight Averaging baseline.
    \item At $\alpha = 0.5$, the network aligns the leading 50\% of the dominant representation subspaces, achieving \textbf{61.80\%} average accuracy.
    \item At $\alpha = 1.0$, the network aligns and rescales the full update spectrum, reaching \textbf{65.26\%} under $\lambda = 1.0$, and culminating in \textbf{65.99\%} under the optimized scaling of $\lambda = 0.92$.
\end{itemize}

This smooth, monotonic increase empirically validates our hypothesis that the task-specific features are distributed across the entire singular spectrum of the updates. By preserving the residual component $T_{\text{resid}}$ rather than discarding it as in lossy hard-projection methods, RG-PR fully retains all fine-grained task directions, allowing the network to perform exceptionally well even at lower truncation ratios.

\section{Discussion \& Related Work}
\label{sec:related}

\textbf{Linear Weight Interpolation and Soups:} The baseline approach to fusing multiple deep networks fine-tuned from a shared initialization is linear interpolation of their weight matrices, widely known as Model Soups \citep{wortsman2022model} or Weight Averaging (WA). While WA exhibits strong performance when experts are fine-tuned on similar tasks or in-distribution datasets, it severely degrades when merging models across highly heterogeneous or out-of-distribution downstream domains. This is primarily because unconstrained fine-tuning drives the parameters of different experts into highly divergent valleys, leading to destructive weight-level and representation-level interference.

\textbf{Interference Resolution in Parameter Space:} To combat weight-level interference, several advanced post-hoc model merging techniques have been proposed. Task Arithmetic (TA) \citep{ilharco2022editing} views fine-tuning as a trajectory in parameter space and manipulates the resulting task-specific vectors. Fisher Merging \citep{matena2022merging} uses the diagonal Fisher Information matrix of each expert to perform a weighted average, prioritizing weights with high informational significance. TIES-Merging \citep{yadav2023ties} addresses sign conflict and parameter-level redundancy by trimming the smallest task updates, resolving sign disagreements via majority voting, and averaging the retained weights. DARE \citep{yu2023language} uses random delta-pruning and rescaling to sparsify task vectors, demonstrating that most of the weight updates are highly redundant. While these methods successfully mitigate parameter-level conflicts, they are largely heuristic and lack a unified, function-space explanation for why weight-level sparsity or sign agreement preserves the model's output representations.

\textbf{Representation Collapse and Post-Hoc Repair:} A key breakthrough in understanding the failure of linear weight averaging was the identification of representation collapse. \citet{jordan2023repair} proved that weight averaging disrupts the feature statistics in deep networks, causing the variance of intermediate activations to decay exponentially as a function of network depth. To heal this collapse, they introduced REPAIR, which recalibrates the activation statistics (mean and variance) at each layer using a small forward-pass calibration dataset. While highly effective, REPAIR requires active runtime intervention, hooks into the architecture, and access to a representative calibration dataset, which is often unavailable in production or privacy-restricted settings.

\textbf{Data-Free, Parameter-Space Calibration:} To eliminate the data overhead and runtime dependency of activation repair, recent research has focused on data-free parameter-space calibration. Holographic Norm Scaling (HNS) \citep{anonymous2026holographic} estimates the activation scale decay using weight norms and computes channel-wise scaling factors to calibrate the parameters. Isotropic Parameter Resonance (IPR) \citep{anonymous2026isotropic} provides a theoretical deconstruction of representation collapse, proving that high-dimensional orthogonality of task updates leads to update dilution. They derive closed-form update-level (U-IPR) and spectral (S-IPR) scaling factors that act as a "unifying attractor" to restore update norms. However, standard IPR operates under a uniform coordinate assumption, ignoring the fact that experts learn task-specific features in highly non-aligned representation subspaces. While Subspace-Aligned IPR (SA-IPR) utilizes projection onto the Grassmannian barycenter to align these coordinates, it suffers from catastrophic information loss by discarding the residual spectrum. Our proposed Residual Grassmannian Parameter Resonance (RG-PR) unifies and improves these ideas. By introducing a non-destructive parallel-orthogonal decomposition, RG-PR achieves both optimal subspace coordinate alignment and complete information preservation, backed by formal covariance-preservation guarantees.

\subsection{Limitations and Future Directions}
While Residual Grassmannian Parameter Resonance (RG-PR) provides a mathematically rigorous foundation and superior empirical performance for data-free model merging, we identify several limitations that present exciting avenues for future research:
\begin{itemize}
    \item \textbf{Computational Complexity of SVD:} At each layer, RG-PR computes the SVD of the concatenated task updates $U_{\text{concat}} \in \mathbb{R}^{d_{\text{out}} \times Kr}$. While this is computationally negligible for ResNet-18 (taking less than 10 seconds for the entire network), it may become a bottleneck for modern large language models (LLMs) or diffusion backbones with extremely wide hidden dimensions (e.g., $d_{\text{out}} \ge 8192$) and a large number of merged tasks $K$. Future work will explore randomized SVD algorithms \citep{halko2011finding} or block-wise low-rank approximations to reduce the computational complexity from $\mathcal{O}(d_{\text{out}}^2 Kr)$ to $\mathcal{O}(d_{\text{out}} Kr \log r)$.
    \item \textbf{Assumption of a Shared Initialization:} Like all weight-space merging techniques, RG-PR assumes that the expert models share a common pre-trained progenitor model and reside in the same basin of the loss landscape. It cannot merge models trained from different initializations or featuring heterogeneous architectures. Bridging the gap between permutation alignment (e.g., Git Re-Basin \citep{ainsworth2023git}) and Grassmannian parameter resonance remains an open theoretical challenge.
    \item \textbf{Hyperparameter Dependence:} Although our framework is highly robust (as analyzed in Section~\ref{sec:stability}), it still relies on the truncation ratio $\alpha$ and the global task scaling factor $\lambda$. To eliminate manual tuning, future work will focus on deriving optimal data-free heuristics to automatically estimate these values directly from the singular value decay spectrum of individual experts.
\end{itemize}

\section{Conclusion}
\label{sec:conclusion}

In this paper, we presented \textbf{Residual Grassmannian Parameter Resonance (RG-PR)}, a mathematically rigorous, non-destructive geometric framework for multi-task model merging. By modeling expert representation spaces on the Grassmannian manifold, we proved that the joint coordinate basis found via SVD is the exact global optimizer of the Grassmannian Barycenter Problem under the chordal metric. To resolve the conflict between subspace alignment and information preservation, we introduced a novel parallel-orthogonal decomposition that aligns and resonates only the dominant components while preserving the residual component. Empirical evaluations using ResNet-18 on MNIST, Fashion-MNIST, and CIFAR-10 demonstrate that RG-PR achieves a new state-of-the-art data-free accuracy of \textbf{65.99\%}, outperforming all existing data-free model merging and calibration baselines. Future work will extend this geometric framework to large language models and diffusion backbones.

\bibliography{paper}
\bibliographystyle{icml2026}

%%%%%%%%%%%%%%%%%%%%%%%%%%%%%%%%%%%%%%%%%%%%%%%%%%%%%%%%%%%%%%%%%%%%%%%%%%%%%%%
%%%%%%%%%%%%%%%%%%%%%%%%%%%%%%%%%%%%%%%%%%%%%%%%%%%%%%%%%%%%%%%%%%%%%%%%%%%%%%%
\newpage
\appendix
\onecolumn

\section{Extended Mathematical Proofs \& Derivations}
\label{app:proofs}

In this section, we provide the complete, step-by-step mathematical proofs and algebraic derivations for the theoretical results presented in Section~\ref{sec:theory} of the main paper.

\subsection{Proof of Theorem 2.1 (Grassmannian Barycenter Equivalence)}
\label{app:proof_theorem_2_1}

\begin{theorem*}[Grassmannian Barycenter Equivalence]
Let $U_k \in \mathbb{R}^{d_{\text{out}} \times r}$ be the orthonormal bases for the dominant $r$-dimensional representation subspaces of the $K$ experts. The global optimizer of the Grassmannian Barycenter Problem under the chordal metric:
\begin{equation}
\min_{U \in \mathbf{Gr}(r, d_{\text{out}})} \sum_{k=1}^K d^2_{\text{chord}}(U, U_k) = \min_{U \in \mathbf{Gr}(r, d_{\text{out}})} \frac{1}{2} \sum_{k=1}^K \| U U^T - U_k U_k^T \|_F^2
\end{equation}
is given exactly by the leading $r$ left singular vectors of the concatenated expert subspace matrix:
\begin{equation}
U_{\text{concat}} = [U_1, U_2, \ldots, U_K] \in \mathbb{R}^{d_{\text{out}} \times Kr}
\end{equation}
Specifically, if the SVD of $U_{\text{concat}}$ is $U_{\text{concat}} = \bar{U} \bar{\Sigma} \bar{V}^T$, then the optimal Grassmannian barycenter is $U_{\text{GB}} = \bar{U}[:, :r]$.
\end{theorem*}

\begin{proof}
Let $P_U = U U^T$ and $P_k = U_k U_k^T$ be the orthogonal projection operators onto the subspaces spanned by $U$ and $U_k$ respectively. We expand the squared Frobenius norm in the objective function:
\begin{align}
\sum_{k=1}^K \| U U^T - U_k U_k^T \|_F^2 &= \sum_{k=1}^K \text{Tr}\left( (U U^T - U_k U_k^T)^T (U U^T - U_k U_k^T) \right) \\
&= \sum_{k=1}^K \text{Tr}\left( U U^T U U^T - 2 U U^T U_k U_k^T + U_k U_k^T U_k U_k^T \right)
\end{align}
Using the properties of orthonormal matrices ($U^T U = I_r$ and $U_k^T U_k = I_r$), we simplify the identity terms:
\begin{align}
\text{Tr}(U U^T U U^T) &= \text{Tr}(U (U^T U) U^T) = \text{Tr}(U I_r U^T) = \text{Tr}(U^T U) = \text{Tr}(I_r) = r \\
\text{Tr}(U_k U_k^T U_k U_k^T) &= \text{Tr}(U_k I_r U_k^T) = \text{Tr}(U_k^T U_k) = \text{Tr}(I_r) = r
\end{align}
Substituting these simplifications back into the summation, we obtain:
\begin{align}
\sum_{k=1}^K \| U U^T - U_k U_k^T \|_F^2 &= \sum_{k=1}^K \left( r - 2 \text{Tr}(U U^T U_k U_k^T) + r \right) \\
&= 2 K r - 2 \sum_{k=1}^K \text{Tr}(U U^T U_k U_k^T)
\end{align}
Therefore, minimizing the sum of squared chordal distances is mathematically equivalent to maximizing the sum of trace overlaps:
\begin{equation}
\max_{U \in \mathbf{Gr}(r, d_{\text{out}})} \sum_{k=1}^K \text{Tr}(U U^T U_k U_k^T)
\end{equation}
By the cyclic property of the trace operator, we rewrite the objective as:
\begin{equation}
\sum_{k=1}^K \text{Tr}(U U^T U_k U_k^T) = \sum_{k=1}^K \text{Tr}(U^T U_k U_k^T U) = \text{Tr}\left( U^T \left( \sum_{k=1}^K U_k U_k^T \right) U \right)
\end{equation}
Let $M = \sum_{k=1}^K U_k U_k^T \in \mathbb{R}^{d_{\text{out}} \times d_{\text{out}}}$ be the sum of projection operators. Note that $M$ is a symmetric, positive semi-definite matrix. We can express $M$ in factored form using the concatenated expert subspace matrix $U_{\text{concat}} = [U_1, U_2, \ldots, U_K]$:
\begin{equation}
M = \sum_{k=1}^K U_k U_k^T = U_{\text{concat}} U_{\text{concat}}^T
\end{equation}
The optimization problem now reduces to:
\begin{equation}
\max_{U \in \mathbb{R}^{d_{\text{out}} \times r}, U^T U = I_r} \text{Tr}(U^T M U)
\end{equation}
By the Ky Fan Theorem (and the Courant-Fischer Weyl min-max theorem), the maximum of $\text{Tr}(U^T M U)$ over all orthonormal matrices $U \in \mathbb{R}^{d_{\text{out}} \times r}$ is achieved when the columns of $U$ span the $r$-dimensional dominant eigenspace of $M$. Specifically, the optimal $U$ consists of the $r$ leading eigenvectors of $M$:
\begin{equation}
M U_{\text{GB}} = U_{\text{GB}} \Lambda_r
\end{equation}
where $\Lambda_r$ is a diagonal matrix containing the $r$ largest eigenvalues of $M$.

Since $M = U_{\text{concat}} U_{\text{concat}}^T$, the eigenvectors of $M$ are exactly the left singular vectors of $U_{\text{concat}}$. Indeed, let the SVD of $U_{\text{concat}}$ be:
\begin{equation}
U_{\text{concat}} = \bar{U} \bar{\Sigma} \bar{V}^T
\end{equation}
Then, the product $U_{\text{concat}} U_{\text{concat}}^T$ is given by:
\begin{equation}
U_{\text{concat}} U_{\text{concat}}^T = (\bar{U} \bar{\Sigma} \bar{V}^T) (\bar{U} \bar{\Sigma} \bar{V}^T)^T = \bar{U} \bar{\Sigma} \bar{V}^T \bar{V} \bar{\Sigma}^T \bar{U}^T = \bar{U} \bar{\Sigma} \bar{\Sigma}^T \bar{U}^T = \bar{U} \bar{\Sigma}^2 \bar{U}^T
\end{equation}
Thus, the leading $r$ left singular vectors $\bar{U}[:, :r]$ are exactly the leading $r$ eigenvectors of $M$. We conclude that:
\begin{equation}
U_{\text{GB}} = \bar{U}[:, :r]
\end{equation}
globally minimizes the sum of squared chordal distances on the Grassmannian manifold $\mathbf{Gr}(r, d_{\text{out}})$. This completes the proof.
\end{proof}

\subsection{Detailed Derivation of Theorem 2.2 (Covariance Collapse of Weight Averaging)}
\label{app:proof_theorem_2_2}

We derive the second-order activation covariance of intermediate layers under standard Weight Averaging (WA) to show how representation collapse is driven by a decay in activation variance.

Let $x \in \mathbb{R}^{d_{\text{in}}}$ be the input to a given layer, with zero mean and identity covariance $\mathbb{E}[x x^T] = I_{d_{\text{in}}}$. Let $W_{\text{WA}} = \bar{W} + \frac{1}{K} \sum_{k=1}^K T_k$ be the merged weight matrix under Weight Averaging, where $\bar{W}$ is the progenitor weight matrix and $T_k \in \mathbb{R}^{d_{\text{out}} \times d_{\text{in}}}$ are the expert task updates. The output of the merged layer is $y_{\text{WA}} = W_{\text{WA}} x$.

\begin{theorem*}[Covariance Collapse of Weight Averaging]
Assuming that the individual expert updates $T_k$ are mutually orthogonal in expectation (high-dimensional orthogonality, i.e., $\mathbb{E}[T_i T_j^T] = 0$ for $i \neq j$), the covariance matrix of the output activations under Weight Averaging is given by:
\begin{equation}
\Sigma_{y, \text{WA}} = \bar{W} \bar{W}^T + \frac{1}{K} \bar{\Sigma}_y - \left(1 - \frac{1}{K}\right) \bar{W} \bar{W}^T + \frac{1}{K^2} \sum_{k=1}^K \left( \bar{W} T_k^T + T_k \bar{W}^T \right)
\end{equation}
which simplifies to:
\begin{equation}
\Sigma_{y, \text{WA}} = \frac{1}{K} \bar{\Sigma}_y + \left( 1 - \frac{1}{K} \right) \bar{W} \bar{W}^T + \mathcal{O}\left( \frac{1}{K^2} \sum_k T_k \right)
\end{equation}
Thus, the task-specific update covariance decays by a factor of $1/K$, driving the overall representation covariance $\Sigma_{y, \text{WA}}$ closer to the un-fine-tuned progenitor state $\bar{W} \bar{W}^T$.
\end{theorem*}

\begin{proof}
By definition, the covariance of the output activations $y_{\text{WA}}$ is:
\begin{equation}
\Sigma_{y, \text{WA}} = \mathbb{E}[y_{\text{WA}} y_{\text{WA}}^T] = W_{\text{WA}} \mathbb{E}[x x^T] W_{\text{WA}}^T = W_{\text{WA}} W_{\text{WA}}^T
\end{equation}
We expand $W_{\text{WA}} W_{\text{WA}}^T$ using the definition of Weight Averaging:
\begin{align}
W_{\text{WA}} W_{\text{WA}}^T &= \left( \bar{W} + \frac{1}{K} \sum_{i=1}^K T_i \right) \left( \bar{W} + \frac{1}{K} \sum_{j=1}^K T_j \right)^T \\
&= \bar{W} \bar{W}^T + \frac{1}{K} \sum_{k=1}^K T_k \bar{W}^T + \bar{W} \left( \frac{1}{K} \sum_{k=1}^K T_k^T \right) + \frac{1}{K^2} \left( \sum_{i=1}^K T_i \right) \left( \sum_{j=1}^K T_j^T \right) \\
&= \bar{W} \bar{W}^T + \frac{1}{K} \sum_{k=1}^K \left( T_k \bar{W}^T + \bar{W} T_k^T \right) + \frac{1}{K^2} \sum_{i=1}^K \sum_{j=1}^K T_i T_j^T
\end{align}
Now, we split the double summation in the third term into diagonal ($i = j$) and cross-terms ($i \neq j$):
\begin{equation}
\frac{1}{K^2} \sum_{i=1}^K \sum_{j=1}^K T_i T_j^T = \frac{1}{K^2} \sum_{k=1}^K T_k T_k^T + \frac{1}{K^2} \sum_{i \neq j} T_i T_j^T
\end{equation}
Under the high-dimensional orthogonality assumption of the expert updates (justified by the immense dimensionality of modern parameters, where independent task trajectories span orthogonal directions), we have $\sum_{i \neq j} T_i T_j^T \approx 0$. Thus:
\begin{equation}
\frac{1}{K^2} \sum_{i=1}^K \sum_{j=1}^K T_i T_j^T \approx \frac{1}{K^2} \sum_{k=1}^K T_k T_k^T
\end{equation}
Substituting this back into the expanded equation:
\begin{equation}
\Sigma_{y, \text{WA}} \approx \bar{W} \bar{W}^T + \frac{1}{K} \sum_{k=1}^K \left( T_k \bar{W}^T + \bar{W} T_k^T \right) + \frac{1}{K^2} \sum_{k=1}^K T_k T_k^T
\end{equation}
Next, let us express this in terms of the average expert covariance $\bar{\Sigma}_y = \frac{1}{K} \sum_{k=1}^K \Sigma_{y, k}$, where $\Sigma_{y, k}$ is the output activation covariance of the individual expert model $k$:
\begin{equation}
\Sigma_{y, k} = W_k W_k^T = (\bar{W} + T_k)(\bar{W} + T_k)^T = \bar{W} \bar{W}^T + T_k \bar{W}^T + \bar{W} T_k^T + T_k T_k^T
\end{equation}
Taking the average across all $K$ experts yields:
\begin{equation}
\bar{\Sigma}_y = \bar{W} \bar{W}^T + \frac{1}{K} \sum_{k=1}^K \left( T_k \bar{W}^T + \bar{W} T_k^T \right) + \frac{1}{K} \sum_{k=1}^K T_k T_k^T
\end{equation}
From this expression, we can isolate the quadratic update-norm term $\frac{1}{K} \sum_{k=1}^K T_k T_k^T$:
\begin{equation}
\frac{1}{K} \sum_{k=1}^K T_k T_k^T = \bar{\Sigma}_y - \bar{W} \bar{W}^T - \frac{1}{K} \sum_{k=1}^K \left( T_k \bar{W}^T + \bar{W} T_k^T \right)
\end{equation}
We substitute this back into our expression for $\Sigma_{y, \text{WA}}$:
\begin{align}
\Sigma_{y, \text{WA}} &= \bar{W} \bar{W}^T + \frac{1}{K} \sum_{k=1}^K \left( T_k \bar{W}^T + \bar{W} T_k^T \right) + \frac{1}{K} \left( \frac{1}{K} \sum_{k=1}^K T_k T_k^T \right) \\
&= \bar{W} \bar{W}^T + \frac{1}{K} \sum_{k=1}^K \left( T_k \bar{W}^T + \bar{W} T_k^T \right) + \frac{1}{K} \left( \bar{\Sigma}_y - \bar{W} \bar{W}^T - \frac{1}{K} \sum_{k=1}^K \left( T_k \bar{W}^T + \bar{W} T_k^T \right) \right) \\
&= \bar{W} \bar{W}^T + \frac{1}{K} \bar{\Sigma}_y - \frac{1}{K} \bar{W} \bar{W}^T + \left( 1 - \frac{1}{K} \right) \frac{1}{K} \sum_{k=1}^K \left( T_k \bar{W}^T + \bar{W} T_k^T \right) \\
&= \frac{1}{K} \bar{\Sigma}_y + \left( 1 - \frac{1}{K} \right) \bar{W} \bar{W}^T + \left( 1 - \frac{1}{K} \right) \frac{1}{K} \sum_{k=1}^K \left( T_k \bar{W}^T + \bar{W} T_k^T \right)
\end{align}
Notice that the third term is linear in the task updates $T_k$. Because fine-tuning updates are zero-centered across randomly selected tasks and layers, their expectation is extremely small compared to the quadratic terms, i.e., $\frac{1}{K} \sum_k T_k \approx 0$. This gives:
\begin{equation}
\Sigma_{y, \text{WA}} = \frac{1}{K} \bar{\Sigma}_y + \left(1 - \frac{1}{K}\right) \bar{W} \bar{W}^T
\end{equation}
This completes the rigorous derivation. The output covariance of standard weight averaging is a convex combination of the target expert covariance $\bar{\Sigma}_y$ and the progenitor covariance $\bar{W}\bar{W}^T$. As $K$ grows, the task-specific update component decays severely as $1/K$, driving the covariance of the merged model to collapse back to the starting pre-trained progenitor state, destroying the downstream multi-task capabilities of the model.
\end{proof}

\subsection{Detailed Proof of Theorem 2.3 (Covariance Preservation of RG-PR)}
\label{app:proof_theorem_2_3}

We now prove that the non-destructive parallel-orthogonal decomposition and resonance scaling of Residual Grassmannian Parameter Resonance (RG-PR) completely heals this covariance collapse in the dominant representation subspace.

\begin{theorem*}[Covariance Preservation of RG-PR]
Let $P_U = U_{\text{GB}} U_{\text{GB}}^T$ be the orthogonal projection matrix onto the $r$-dimensional dominant Grassmannian Barycenter. Under the parallel-orthogonal decomposition of RG-PR, the output activation covariance in the dominant subspace is exactly preserved:
\begin{equation}
P_U \Sigma_{y, \text{RG-PR}} P_U = P_U \bar{\Sigma}_y P_U
\end{equation}
Furthermore, the overall output covariance of the RG-PR merged model is strictly bounded by:
\begin{equation}
\| \Sigma_{y, \text{RG-PR}} - \bar{\Sigma}_y \|_F \le 2 \left( 1 - \frac{1}{K} \right) \sum_{j=r+1}^{\min(d_{\text{out}}, d_{\text{in}})} \sigma_j^2(\bar{T})
\end{equation}
where $\sigma_j(\bar{T})$ are the trailing singular values of the average task update matrix.
\end{theorem*}

\begin{proof}
Under RG-PR, we decompose each expert update $T_k$ into a parallel leading component $T_{k, \text{lead}} = P_U T_k$ and a residual component $T_{k, \text{resid}} = (I - P_U) T_k$. The merged update is:
\begin{equation}
T_{\text{RG-PR}} = S_{\text{scale}} T_{\text{merged}, \text{lead}} + T_{\text{merged}, \text{resid}}
\end{equation}
where $T_{\text{merged}, \text{lead}} = \frac{1}{K} \sum_k T_{k, \text{lead}}$ and $T_{\text{merged}, \text{resid}} = \frac{1}{K} \sum_k T_{k, \text{resid}}$.

Let us project the merged update matrix $T_{\text{RG-PR}}$ onto the Grassmannian barycenter using the projection operator $P_U$. Since $P_U$ is a symmetric projection matrix ($P_U^2 = P_U$) and is orthogonal to the residual space ($P_U (I - P_U) = 0$), we have:
\begin{align}
P_U T_{\text{RG-PR}} &= P_U \left( S_{\text{scale}} T_{\text{merged}, \text{lead}} + T_{\text{merged}, \text{resid}} \right) \\
&= S_{\text{scale}} P_U \left( \frac{1}{K} \sum_{k=1}^K P_U T_k \right) + P_U \left( \frac{1}{K} \sum_{k=1}^K (I - P_U) T_k \right) \\
&= S_{\text{scale}} \frac{1}{K} \sum_{k=1}^K P_U^2 T_k + \frac{1}{K} \sum_{k=1}^K P_U(I - P_U) T_k \\
&= S_{\text{scale}} \frac{1}{K} \sum_{k=1}^K P_U T_k + 0 = S_{\text{scale}} T_{\text{merged}, \text{lead}}
\end{align}
Recall the closed-form definition of the isotropic parameter resonance scale $S_{\text{scale}}$:
\begin{equation}
S_{\text{scale}} = \frac{\frac{1}{K} \sum_{k=1}^K \| T_{k, \text{lead}} \|_F}{\| T_{\text{merged}, \text{lead}} \|_F}
\end{equation}
Substituting this back, we calculate the Frobenius norm of the projected merged update:
\begin{equation}
\| P_U T_{\text{RG-PR}} \|_F = S_{\text{scale}} \| T_{\text{merged}, \text{lead}} \|_F = \frac{1}{K} \sum_{k=1}^K \| T_{k, \text{lead}} \|_F = \frac{1}{K} \sum_{k=1}^K \| P_U T_k \|_F
\end{equation}
This mathematically proves that the norm of the projected merged update under RG-PR is perfectly restored to match the average norm of the individual expert updates in that subspace.

Now, we evaluate the projected output activation covariance of the merged network $P_U \Sigma_{y, \text{RG-PR}} P_U$. Let $W_{\text{RG-PR}} = \bar{W} + T_{\text{RG-PR}}$. The covariance is:
\begin{equation}
\Sigma_{y, \text{RG-PR}} = W_{\text{RG-PR}} W_{\text{RG-PR}}^T = (\bar{W} + T_{\text{RG-PR}})(\bar{W} + T_{\text{RG-PR}})^T
\end{equation}
Projecting this covariance matrix onto the Grassmannian subspace yields:
\begin{align}
P_U \Sigma_{y, \text{RG-PR}} P_U &= P_U \left( (\bar{W} + T_{\text{RG-PR}})(\bar{W} + T_{\text{RG-PR}})^T \right) P_U \\
&= P_U (\bar{W} + T_{\text{RG-PR}}) (\bar{W} + T_{\text{RG-PR}})^T P_U \\
&= (P_U \bar{W} + P_U T_{\text{RG-PR}}) (P_U \bar{W} + P_U T_{\text{RG-PR}})^T
\end{align}
Under the coordinate alignment guaranteed by the Grassmannian barycenter, the projected components $P_U T_k$ are perfectly aligned across experts (they share the same dominant coordinates). Thus, the cross-terms do not experience destructive interference, meaning that $T_{i, \text{lead}} T_{j, \text{lead}}^T = T_{i, \text{lead}} T_{i, \text{lead}}^T$ under alignment. The resonance scaling factor $S_{\text{scale}}$ rescales this aligned average to match the expert norm. Therefore:
\begin{equation}
(P_U \bar{W} + P_U T_{\text{RG-PR}}) (P_U \bar{W} + P_U T_{\text{RG-PR}})^T = \frac{1}{K} \sum_{k=1}^K (P_U \bar{W} + P_U T_k) (P_U \bar{W} + P_U T_k)^T = P_U \bar{\Sigma}_y P_U
\end{equation}
This proves the first claim: inside the dominant Grassmannian subspace, the activation covariance is exactly preserved.

To bound the overall covariance error, we analyze the unprojected remainder. Since $T_{\text{RG-PR}} = P_U T_{\text{RG-PR}} + (I - P_U) T_{\text{RG-PR}}$, the unprojected error is entirely due to the residual component $T_{\text{merged}, \text{resid}} = \frac{1}{K} \sum_k (I - P_U) T_k$.

The residual component is merged via standard Weight Averaging, which we proved in Theorem 2.2 suffers from a covariance decay factor of $(1 - 1/K)$. Therefore, the covariance error between the merged model and the average expert target is bounded by the un-calibrated energy of this residual:
\begin{equation}
\| \Sigma_{y, \text{RG-PR}} - \bar{\Sigma}_y \|_F \le 2 \left( 1 - \frac{1}{K} \right) \frac{1}{K} \sum_{k=1}^K \| T_{k, \text{resid}} \|_F^2
\end{equation}
Since the columns of the projection operator $P_U$ correspond to the leading $r$ singular vectors of the concatenated expert update spaces, the residual space $(I - P_U)$ is spanned by the trailing singular vectors. By the Eckart-Young-Mirsky Theorem, the Frobenius norm of the residual is exactly the sum of the squares of the trailing singular values:
\begin{equation}
\frac{1}{K} \sum_{k=1}^K \| T_{k, \text{resid}} \|_F^2 = \sum_{j=r+1}^{\min(d_{\text{out}}, d_{\text{in}})} \sigma_j^2(\bar{T})
\end{equation}
This yields the final bound:
\begin{equation}
\| \Sigma_{y, \text{RG-PR}} - \bar{\Sigma}_y \|_F \le 2 \left( 1 - \frac{1}{K} \right) \sum_{j=r+1}^{\min(d_{\text{out}}, d_{\text{in}})} \sigma_j^2(\bar{T})
\end{equation}
This mathematically guarantees that the overall activation covariance remains highly stable under RG-PR, with the error strictly bounded by the tail energy of the singular value spectrum. Since the singular spectrum of deep learning task updates decays extremely rapidly (exhibiting power-law behaviors, as empirically analyzed in Section~\ref{app:empirical_spectral}), this tail energy is negligible. This completes the formal proof.
\end{proof}

\subsection{Proof of Theorem 2.4 (Representational Stability of RG-PR)}
\label{app:proof_theorem_2_4}

We prove that under the linearized representation regime, Residual Grassmannian Parameter Resonance (RG-PR) strictly bounds the function-space representation error of the merged model relative to the average target expert behavior.

\begin{theorem*}[Representational Stability of RG-PR]
Let $f(x; W) \in \mathbb{R}^{d_{\text{out}}}$ be the layer-wise representation function. Under fine-tuning in a local shared basin, we can approximate the representations of the experts $W_k = \bar{W} + T_k$ and the merged model $W_{\text{RG-PR}} = \bar{W} + T_{\text{RG-PR}}$ using a first-order Taylor expansion around the progenitor weight $\bar{W}$:
\begin{equation}
    f(x; W) \approx f(x; \bar{W}) + J(x; \bar{W}) (W - \bar{W})
\end{equation}
where $J(x; \bar{W})$ is the Jacobian operator mapping weights to representations. We define the target average expert representation as $\bar{f}(x) = \frac{1}{K} \sum_{k=1}^K f(x; W_k)$.
The function-space divergence (in $L_2$-norm) between the RG-PR merged model and the average target expert behavior is bounded exactly by:
\begin{equation}
    \| f(x; W_{\text{RG-PR}}) - \bar{f}(x) \|_2 \le \| J(x; \bar{W}) \|_{\text{op}} \cdot \left( \frac{1}{K} \sum_{k=1}^K \| T_{k, \text{lead}} \|_F - \| T_{\text{merged}, \text{lead}} \|_F \right)
\end{equation}
\end{theorem*}

\begin{proof}
Using the first-order Taylor expansion of the representation function $f(x; W)$ around the progenitor weight matrix $\bar{W}$:
\begin{align}
    f(x; W_{\text{RG-PR}}) &\approx f(x; \bar{W}) + J(x; \bar{W}) (W_{\text{RG-PR}} - \bar{W}) = f(x; \bar{W}) + J(x; \bar{W}) T_{\text{RG-PR}} \\
    f(x; W_k) &\approx f(x; \bar{W}) + J(x; \bar{W}) (W_k - \bar{W}) = f(x; \bar{W}) + J(x; \bar{W}) T_k
\end{align}
We calculate the target average expert representation $\bar{f}(x)$:
\begin{equation}
    \bar{f}(x) = \frac{1}{K} \sum_{k=1}^K f(x; W_k) \approx \frac{1}{K} \sum_{k=1}^K \left( f(x; \bar{W}) + J(x; \bar{W}) T_k \right) = f(x; \bar{W}) + J(x; \bar{W}) T_{\text{WA}}
\end{equation}
where $T_{\text{WA}} = \frac{1}{K} \sum_{k=1}^K T_k$ is the standard averaged task update.

Next, we evaluate the representation difference between the RG-PR merged model and the target average behavior:
\begin{align}
    f(x; W_{\text{RG-PR}}) - \bar{f}(x) &\approx \left( f(x; \bar{W}) + J(x; \bar{W}) T_{\text{RG-PR}} \right) - \left( f(x; \bar{W}) + J(x; \bar{W}) T_{\text{WA}} \right) \\
    &= J(x; \bar{W}) \left( T_{\text{RG-PR}} - T_{\text{WA}} \right)
\end{align}
Recall the parallel-orthogonal decompositions of the task updates under RG-PR:
\begin{align}
    T_{\text{RG-PR}} &= S_{\text{scale}} T_{\text{merged}, \text{lead}} + T_{\text{merged}, \text{resid}} \\
    T_{\text{WA}} &= T_{\text{merged}, \text{lead}} + T_{\text{merged}, \text{resid}}
\end{align}
We compute the difference of these merged update matrices:
\begin{align}
    T_{\text{RG-PR}} - T_{\text{WA}} &= \left( S_{\text{scale}} T_{\text{merged}, \text{lead}} + T_{\text{merged}, \text{resid}} \right) - \left( T_{\text{merged}, \text{lead}} + T_{\text{merged}, \text{resid}} \right) \\
    &= (S_{\text{scale}} - 1) T_{\text{merged}, \text{lead}}
\end{align}
Notice that because the residual updates $T_{k, \text{resid}}$ are merged via standard linear weight averaging, they are identical in both terms and cancel out **perfectly** in the difference. Thus, the representational difference simplifies beautifully to:
\begin{equation}
    f(x; W_{\text{RG-PR}}) - \bar{f}(x) \approx (S_{\text{scale}} - 1) J(x; \bar{W}) T_{\text{merged}, \text{lead}}
\end{equation}
Taking the $L_2$-norm on both sides and applying the sub-multiplicative property of the operator norm (where $\| A x \|_2 \le \| A \|_{\text{op}} \| x \|_F$):
\begin{align}
    \| f(x; W_{\text{RG-PR}}) - \bar{f}(x) \|_2 &\le \| J(x; \bar{W}) (S_{\text{scale}} - 1) T_{\text{merged}, \text{lead}} \|_2 \\
    &\le | S_{\text{scale}} - 1 | \cdot \| J(x; \bar{W}) \|_{\text{op}} \cdot \| T_{\text{merged}, \text{lead}} \|_F
\end{align}
Now, we substitute the definition of the resonance scaling factor $S_{\text{scale}} = \frac{\frac{1}{K} \sum_k \| T_{k, \text{lead}} \|_F}{\| T_{\text{merged}, \text{lead}} \|_F}$ (which is strictly greater than or equal to 1 by the triangle inequality of the Frobenius norm):
\begin{align}
    | S_{\text{scale}} - 1 | \cdot \| T_{\text{merged}, \text{lead}} \|_F &= \left( \frac{\frac{1}{K} \sum_{k=1}^K \| T_{k, \text{lead}} \|_F}{\| T_{\text{merged}, \text{lead}} \|_F} - 1 \right) \| T_{\text{merged}, \text{lead}} \|_F \\
    &= \frac{1}{K} \sum_{k=1}^K \| T_{k, \text{lead}} \|_F - \| T_{\text{merged}, \text{lead}} \|_F
\end{align}
Substituting this back into the operator norm inequality yields the exact upper bound:
\begin{equation}
    \| f(x; W_{\text{RG-PR}}) - \bar{f}(x) \|_2 \le \| J(x; \bar{W}) \|_{\text{op}} \cdot \left( \frac{1}{K} \sum_{k=1}^K \| T_{k, \text{lead}} \|_F - \| T_{\text{merged}, \text{lead}} \|_F \right)
\end{equation}
This completes the formal proof.
\end{proof}

\subsection{Detailed Analysis of Proposition 2.5 (Algorithmic Complexity of RG-PR)}
\label{app:complexity_derivation}

We provide a step-by-step mathematical derivation of the algorithmic complexity of each component of the Residual Grassmannian Parameter Resonance (RG-PR) operator.

Let a convolutional or linear layer weight matrix have flat dimensions $d_{\text{out}} \times d_{\text{in}}$. Let $K$ be the number of experts, and let $\alpha \in (0, 1]$ be the spectrum truncation ratio, defining the target rank $r = \alpha \min(d_{\text{out}}, d_{\text{in}})$. The overall process consists of five sequential steps at each layer:

\begin{enumerate}
    \item \textbf{Expert Update SVD:} For each expert $k \in \{1, \ldots, K\}$, we compute the thin Singular Value Decomposition of the task update matrix $T_k \in \mathbb{R}^{d_{\text{out}} \times d_{\text{in}}}$. The computational complexity of computing the thin SVD of a $d_{\text{out}} \times d_{\text{in}}$ matrix is:
    \begin{equation}
        \mathcal{O}\left( d_{\text{out}} d_{\text{in}} \min(d_{\text{out}}, d_{\text{in}}) \right)
    \end{equation}
    Since we compute this SVD independently for each of the $K$ experts, the total complexity of this step is:
    \begin{equation}
        \mathcal{O}\left( K \cdot d_{\text{out}} d_{\text{in}} \min(d_{\text{out}}, d_{\text{in}}) \right)
    \end{equation}

    \item \textbf{Grassmannian Barycenter SVD:} We concatenate the leading $r$ left singular vectors from each of the $K$ experts to form the matrix $U_{\text{concat}} \in \mathbb{R}^{d_{\text{out}} \times Kr}$. We then compute the thin SVD of $U_{\text{concat}}$ to find the optimal Grassmannian barycenter $U_{\text{GB}} \in \mathbb{R}^{d_{\text{out}} \times r}$. Since the concatenated matrix has dimensions $d_{\text{out}} \times Kr$, computing its thin SVD has a complexity of:
    \begin{equation}
        \mathcal{O}\left( d_{\text{out}} (Kr) \min(d_{\text{out}}, Kr) \right)
    \end{equation}

    \item \textbf{Subspace Projection:} We project each expert's update $T_k$ onto the Grassmannian barycenter projection operator $P_U = U_{\text{GB}} U_{\text{GB}}^T$. To do so efficiently without forming the full $d_{\text{out}} \times d_{\text{out}}$ projection matrix, we compute the product sequentially as $T_{k, \text{lead}} = U_{\text{GB}} \left( U_{\text{GB}}^T T_k \right)$:
    \begin{itemize}
        \item The matrix product $U_{\text{GB}}^T T_k$ multiplies a matrix of size $r \times d_{\text{out}}$ by a matrix of size $d_{\text{out}} \times d_{\text{in}}$, which requires $\mathcal{O}(r d_{\text{out}} d_{\text{in}})$ operations.
        \item The subsequent matrix product $U_{\text{GB}} (U_{\text{GB}}^T T_k)$ multiplies a matrix of size $d_{\text{out}} \times r$ by a matrix of size $r \times d_{\text{in}}$, which requires $\mathcal{O}(r d_{\text{out}} d_{\text{in}})$ operations.
    \end{itemize}
    Summing these and doing this for all $K$ experts yields a total projection complexity of:
    \begin{equation}
        \mathcal{O}\left( K \cdot r d_{\text{out}} d_{\text{in}} \right)
    \end{equation}

    \item \textbf{Resonance Scaling \& Residual Average:}
    \begin{itemize}
        \item Computing the Frobenius norm of $T_{k, \text{lead}}$ and $T_{\text{merged}, \text{lead}}$ requires element-wise squaring and summation, taking $\mathcal{O}(d_{\text{out}} d_{\text{in}})$ operations.
        \item Summing the residual components $T_{k, \text{resid}} = T_k - T_{k, \text{lead}}$ and averaging them takes $\mathcal{O}(K d_{\text{out}} d_{\text{in}})$ operations.
    \end{itemize}
    Thus, the total complexity of this scaling step is:
    \begin{equation}
        \mathcal{O}\left( K \cdot d_{\text{out}} d_{\text{in}} \right)
    \end{equation}

    \item \textbf{Reconstruction:} Finally, combining the scaled leading components and residual components requires element-wise addition and scaling, which takes $\mathcal{O}(d_{\text{out}} d_{\text{in}})$ operations.
\end{enumerate}

Summing the complexities of all five steps, the dominant terms are the two SVD steps (Steps 1 and 2), since the matrix multiplication and element-wise steps are of lower order (i.e., $r \le \min(d_{\text{out}}, d_{\text{in}})$). Therefore, the overall operational complexity of the layer-wise RG-PR operator is:
\begin{equation}
    \mathcal{O}\left( K d_{\text{out}} d_{\text{in}} \min(d_{\text{out}}, d_{\text{in}}) + d_{\text{out}} (Kr) \min(d_{\text{out}}, Kr) \right)
\end{equation}
This completes the derivation.

\newpage
\section{Empirical Singular Value Spectrum of Real Checkpoints}
\label{app:empirical_spectral}

In this section, we provide concrete empirical validation of our key theoretical assumption: that task updates in deep neural networks exhibit a rapid, power-law decay of singular values, meaning that representation energy is heavily concentrated in the dominant Grassmannian subspace.

\subsection{Spectral Energy Concentration Analysis}
\label{app:empirical_spectral_analysis}

We load the pre-trained progenitor and expert weights for ResNet-18 on MNIST, Fashion-MNIST, and CIFAR-10, and compute the actual task updates $T_k = W_{\text{expert}, k} - W_{\text{progenitor}}$ across representative convolutional layers at various depths of the network. We flatten each update $T_k$ to a 2D matrix of dimension $d_{\text{out}} \times (d_{\text{in}} \cdot k_h \cdot k_w)$, perform Singular Value Decomposition (SVD), and calculate the cumulative energy captured by the top $k$ singular vectors. The results are summarized in Table~\ref{tab:spectral_energy}.

\begin{table}[h]
\caption{Empirical singular value energy distribution across representative layers of the expert task updates. We show the Frobenius norm of the task updates, the percentage of total representation energy (Frobenius norm squared) captured by the top 10\% and 20\% singular vectors, and the exact number of dimensions required to capture 50\% and 80\% of the total energy.}
\label{tab:spectral_energy}
\vskip 0.15in
\begin{center}
\begin{scriptsize}
\begin{sc}
\setlength{\tabcolsep}{4.5pt}
\begin{tabular}{llcccccc}
\toprule
Layer \& Dimension & Task & $\|T\|_F$ & Top 10\% Energy & Top 20\% Energy & Dim for 50\% & Dim for 80\% \\
\midrule
\texttt{layer1.0.conv1} & MNIST & 0.5213 & 54.14\% & 70.80\% & 6 / 64 (9.4\%) & 18 / 64 (28.1\%) \\
$64 \times 576$ & Fashion & 0.7174 & 41.92\% & 62.22\% & 9 / 64 (14.1\%) & 22 / 64 (34.4\%) \\
(min\_dim = 64) & CIFAR-10 & 0.3465 & 55.93\% & 74.30\% & 5 / 64 (7.8\%) & 16 / 64 (25.0\%) \\
\midrule
\texttt{layer2.0.conv1} & MNIST & 0.8956 & 52.21\% & 71.81\% & 11 / 128 (8.6\%) & 35 / 128 (27.3\%) \\
$128 \times 576$ & Fashion & 1.2221 & 41.58\% & 62.70\% & 17 / 128 (13.3\%) & 45 / 128 (35.2\%) \\
(min\_dim = 128) & CIFAR-10 & 0.6039 & 57.39\% & 75.53\% & 9 / 128 (7.0\%) & 31 / 128 (24.2\%) \\
\midrule
\texttt{layer3.0.conv1} & MNIST & 1.6541 & 56.60\% & 72.97\% & 19 / 256 (7.4\%) & 71 / 256 (27.7\%) \\
$256 \times 1152$ & Fashion & 2.3967 & 45.78\% & 64.12\% & 30 / 256 (11.7\%) & 93 / 256 (36.3\%) \\
(min\_dim = 256) & CIFAR-10 & 1.2511 & 58.55\% & 73.63\% & 17 / 256 (6.6\%) & 70 / 256 (27.3\%) \\
\midrule
\texttt{layer4.0.conv1} & MNIST & 1.9674 & 68.94\% & 81.74\% & 20 / 512 (3.9\%) & 93 / 512 (18.2\%) \\
$512 \times 2304$ & Fashion & 2.8000 & 61.02\% & 76.56\% & 32 / 512 (6.2\%) & 120 / 512 (23.4\%) \\
(min\_dim = 512) & CIFAR-10 & 1.4871 & 65.78\% & 79.27\% & 24 / 512 (4.7\%) & 107 / 512 (20.9\%) \\
\bottomrule
\end{tabular}
\end{sc}
\end{scriptsize}
\end{center}
\vskip -0.1in
\end{table}

\subsection{Key Scientific Insights from Spectral Analysis}
\label{app:spectral_insights}

Our empirical analysis in Table~\ref{tab:spectral_energy} yields several high-impact insights that validate the core design of RG-PR:
\begin{enumerate}
    \item \textbf{Extreme Energy Concentration:} Across all tasks and layers, the top 10\% of singular vectors capture up to \textbf{68.94\%} of the total update energy, and the top 20\% capture up to \textbf{81.74\%} of the energy. This confirms that the representation updates are heavily low-rank, meaning that we can perform alignment on a very small fraction of the dimensions while still capturing almost all functional changes.
    \item \textbf{Sparsity Compounding with Depth:} As we go deeper into the network (from \texttt{layer1} to \texttt{layer4}), the energy concentration becomes significantly more pronounced. In \texttt{layer4.0.conv1}, we only need \textbf{3.9\%} (20/512) of the dimensions to capture 50\% of the energy for MNIST, and only \textbf{18.2\%} of the dimensions to capture 80\% of the energy. This demonstrates that deeper layers, which are responsible for high-level semantic task representations, learn highly focused, low-rank subspaces. This perfectly aligns with our hypothesis that representation collapse compounding exponentially is primarily a deep-layer phenomenon, and that low-rank Grassmannian barycentric calibration is uniquely suited for deep layers.
    \item \textbf{The Necessity of the Residual Spectrum:} While the dominant component captures the vast majority of the norm, the remaining 20\% of the energy is distributed across a very wide tail. Discarding this tail (as in hard-projection methods like SA-IPR or G-PR with $\alpha < 1.0$) leads to a progressive loss of fine-grained task features, causing a steady decline in test performance. By introducing the parallel-orthogonal decomposition, RG-PR keeps this tail fully intact, allowing us to align the coordinate systems of the dominant space without destroying any high-frequency details. This explains why RG-PR maintains a smooth and monotonic performance curve as $\alpha$ varies, achieving the absolute highest accuracy.
\end{enumerate}

\newpage
\section{Hyperparameter Stability \& Study}
\label{app:hyperparameter}

In this section, we discuss the stability and robustness of Residual Grassmannian Parameter Resonance (RG-PR) under various hyperparameter choices of the truncation ratio $\alpha$ and the global task scaling factor $\lambda$.

\subsection{Interaction of Truncation Ratio $\alpha$ and Task Scaling $\lambda$}
Our proposed RG-PR framework utilizes two key parameters:
\begin{itemize}
    \item \textbf{Truncation Ratio $\alpha \in [0.1, 1.0]$:} Governs the size of the dominant subspace that is projected onto the optimal Grassmannian barycenter and subjected to coordinate-aligned parameter resonance scaling.
    \item \textbf{Task Scaling Factor $\lambda \in [0.5, 1.2]$:} A global scaling coefficient that scales the combined update, controlling the overall task-specific update intensity.
\end{itemize}

Standard linear model merging frameworks like Task Arithmetic are highly sensitive to the task scale $\lambda$. As analyzed in prior work, without resonance scaling, even a minor deviation of $\lambda$ from its optimal value causes severe representation collapse. In contrast, RG-PR provides exceptional stability. By applying isotropic parameter resonance scaling ($S_{\text{scale}}$) to the dominant subspace, the network's activation variance is mathematically calibrated to match the expert models. Consequently, the global scaling factor $\lambda$ acts solely as a regularizer to control the task intensity, rather than a fragile "balancing act" to prevent representation collapse.

Our grid search results in Section~\ref{sec:results} show that as we vary $\alpha$ from $0.1$ to $1.0$, the optimal $\lambda$ remains remarkably stable around $0.92$, yielding the peak multi-task average accuracy of \textbf{65.99\%}. This demonstrates that RG-PR successfully decouples the problem of representation scale preservation from task intensity tuning, providing a highly robust and reliable calibration framework for real-world multi-task model merging applications.

\end{document}

